% Header included in pandoc PDF builds via --include-in-header.
%
% Purpose: map Unicode characters that appear in the canonical .md source
% to LaTeX math-mode equivalents, so the PDF renders correctly even with
% the default Latin Modern fonts (which lack text-mode ≥, ≤, κ, etc.).
%
% Keeping the .md source clean (real Unicode) and translating at build
% time means: (a) the source remains human-readable and grep-able,
% (b) the PDF renders portably without requiring exotic system fonts,
% (c) downstream readers viewing the .md on GitHub/Typora/Obsidian see
% the real characters.

\usepackage{newunicodechar}

% Relational operators
\newunicodechar{≥}{\ensuremath{\geq}}
\newunicodechar{≤}{\ensuremath{\leq}}
\newunicodechar{≈}{\ensuremath{\approx}}
\newunicodechar{≠}{\ensuremath{\neq}}

% Greek letters used in citations and statistics
\newunicodechar{κ}{\ensuremath{\kappa}}
\newunicodechar{α}{\ensuremath{\alpha}}
\newunicodechar{β}{\ensuremath{\beta}}
\newunicodechar{Δ}{\ensuremath{\Delta}}
\newunicodechar{σ}{\ensuremath{\sigma}}

% Punctuation and bullets
\newunicodechar{·}{\ensuremath{\cdot}}
\newunicodechar{×}{\ensuremath{\times}}
\newunicodechar{→}{\ensuremath{\to}}
\newunicodechar{←}{\ensuremath{\leftarrow}}
\newunicodechar{⇒}{\ensuremath{\Rightarrow}}
\newunicodechar{∈}{\ensuremath{\in}}

% Quotation marks (occasionally appear in copy-pasted prose)
\newunicodechar{“}{``}
\newunicodechar{”}{''}
\newunicodechar{‘}{`}
\newunicodechar{’}{'}

% Em-dash and en-dash render fine in lmroman; left for clarity if needed
% \newunicodechar{—}{---}
% \newunicodechar{–}{--}
